% Source: https://artofproblemsolving.com/community/c3h1210476_2007_nationals_target_8
% Stable identity: aops-topic-1210476-q01
Sort by how the four digits are grouped.
\begin{itemize}
\item One $4$-digit number: all $24$ arrangements give different totals ($1234$ to $4321$).
\item $3$-digit $+$ $1$-digit, $\overline{abc}+d$: the total $100a+10b+(c+d)$ depends only on $a$, $b$ (since $c+d=10-a-b$), giving $4\cdot3=12$ totals between $127$ and $433$.
\item $2$-digit $+$ $2$-digit: the total is $10S+(10-S)=9S+10$ where $S$ is the sum of the two tens digits, $S\in\{3,4,5,6,7\}$: totals $37,46,55,64,73$.
\item $2$-digit $+$ $1$ $+$ $1$: the total is $10a+(10-a)=9a+10$ for the tens digit $a$: $19,28,37,46$.
\item $1+2+3+4=10$.
\end{itemize}
The groups have $24+12+5+4+1=46$ totals, and the only repeats are $37$ and $46$ (in both two-digit cases); the ranges of the other groups do not overlap. So there are $\boxed{44}$ different totals.
